\documentclass[12pt]{article}
\usepackage[margin=1in]{geometry}
\usepackage{mathtools, amsfonts, amsthm, amssymb}
\usepackage{xcolor}
\usepackage{graphicx, float}
\usepackage[hybrid]{markdown}
\usepackage{comment}

\renewcommand{\qedsymbol}{\(\blacksquare\)}
\newtheorem{proposition}{Proposition}
\newtheorem*{proposition*}{Proposition}


\title{HW 8}
\author{Jonathan Avalos}
\date{\today}

\begin{document}

\maketitle

\section*{Exercise 1}
Collaborators: None

\noindent
Tools and resources used:
\pagebreak

\section*{Exercise 2}
\begin{itemize}
  \item Points attempted: 2
  \item Self-assessed score: 2.
  \item Time spent: 20 minutes.
  \item Difficulty: 2/10.
  \item Questions/Feedback: None.
\end{itemize}

\pagebreak
\section*{Exercise 3}
\begin{itemize}
  \item Points attempted: 10
  \item Self-assessed score: 10/10.
  \item Time spent: 70 minutes.
  \item Difficulty: 5/10.
  \item Questions/Feedback: None.
\end{itemize}
\begin{comment}
Recall that the first fundamental form \(       \mathrm{I} \colon T_pM \times T_pM \to \mathbb{R}     \) is defined intrinsically as the map \[       \mathrm{I}(a \mathbf{r}_u(p) + b \mathbf{r}_v(p), c \mathbf{r}_u(p) + d \mathbf{r}_v(p)) = acE + (ad+bc)F + bdG,     \] where \[       \begin{align*}         E = \langle \mathbf{r}_u, \mathbf{r}_u \rangle, \quad         F = \langle \mathbf{r}_u, \mathbf{r}_v \rangle, \quad         G = \langle \mathbf{r}_v, \mathbf{r}_v \rangle.       \end{align*}     \]
We typically write this as \[       \mathrm{I} = E du^2 + 2F du\, dv + G dv^2     \] or in the matrix notation (with respect to the basis \(\{ \mathbf{r}_u, \mathbf{r}_v \} \)) \[       \mathrm{I}(\mathbf{w}_1,\mathbf{w}_2) = \langle \mathbf{w}_1,\mathbf{w}_2 \rangle = \mathbf{w}_1^T \begin{pmatrix}       E & F \\ F & G \end{pmatrix} \mathbf{w}_2.     \]


Now the second fundamental form \[       \mathrm{I\!I} \colon T_pM \times T_pM \to \mathbb{R}     \] is defined extrinsically with respect to the Gauss map \(       \mathbf{N} \colon M \to S^2     \) and its differential \( d\mathbf{N}_p \colon T_pM \to T_{\mathbf{N}(p)}(S^2) \), namely: \[       \mathrm{I\!I}(\mathbf{w}_1,\mathbf{w}_2) = \langle -d\mathbf{N}_p(\mathbf{w}_1), \mathbf{w}_2 \rangle = \langle P(\mathbf{w}_1), \mathbf{w}_2 \rangle.     \]
Here, \( P = -d\mathbf{N}_p \) is called the shape operator.1


We will prove that with respect to the basis \(\{ \mathbf{r}_u, \mathbf{r}_v \} \), \[       \mathrm{I\!I}(\mathbf{w}_1,\mathbf{w}_2) = \mathbf{w}_1^T \begin{pmatrix} L & M \\ M & N\end{pmatrix} \mathbf{w}_2,     \] where \[       L = -\mathbf{r}_{uu} \cdot \mathbf{N}, \quad       M = -\mathbf{r}_{uv} \cdot \mathbf{N}, \quad       N = -\mathbf{r}_{vv} \cdot \mathbf{N}.     \]
(Here, \( M \) is a real number, and is different from our surface \( M \).)


Lastly, we will use the first and second fundamental forms (which are easy to compute) to determine the Gaussian and mean curvatures at every point on our surface.
\end{comment}
\begin{enumerate}
  \item 
\end{enumerate}

\end{document}